\chapter{相关技术与理论基础}

\section{MRI 医学图像与数据特点}
MRI 是一种医学影像技术。它利用核磁共振现象获得人体组织信息。与 CT 相比，MRI 在软组织成像方面具有较好对比度。MRI 也不产生电离辐射。因此，MRI 常用于脑部、脊柱、关节和腹部等部位的结构观察。

本文关注脑部 MRI 数据。原始 MRI 数据通常不是一张图片。一个病例通常由多张连续切片组成。这些切片共同构成三维体数据。单张切片只能显示某一层面的灰度分布。连续切片才能表达完整的空间结构。

医学图像与自然图像不同。MRI 图像多为灰度图。图像纹理较少，但边界和灰度层次具有医学含义。医学数据采集成本高。公开数据集规模也常小于自然图像数据集。不同设备和扫描参数会带来分布差异。因此，医学图像生成不能只看视觉效果。系统还需要关注结构合理性、灰度稳定性和结果解释边界。

\section{NIfTI 数据格式与预处理方法}
NIfTI 是神经影像领域常用的数据格式。它支持 \texttt{.nii} 和 \texttt{.nii.gz} 等文件形式。它可以保存体素尺寸、空间方向和仿射坐标等信息\cite{nifti2004}。本文使用的 IXI 数据集提供 NIfTI 格式 MRI 数据。该数据适合用于切片提取、二维生成模型训练和三维序列组织\cite{ixiDataset}。

本文预处理流程服务于 StyleGAN2 训练。系统先读取 NIfTI 体数据。系统再沿轴向提取二维切片。系统保留体数据中间部分，以减少头顶部、颈部和边缘空白区域。然后，系统使用 2\% 到 98\% 百分位截断进行灰度归一化。该步骤把像素值映射到 $[0,255]$。最后，系统统一切片分辨率，并过滤均值、标准差和非零像素比例较低的切片。经过处理后，原始 MRI 体数据被转换为 PNG 切片集合。

\section{生成对抗网络基础}
GAN 由生成器 $G$ 和判别器 $D$ 组成。生成器从随机潜变量 $z$ 出发生成样本。判别器判断输入样本是真实数据还是生成数据。两者通过对抗训练共同优化。经典目标函数可写为：

\begin{equation}
\min_G \max_D V(D,G)=\mathbb{E}_{x\sim p_{data}(x)}[\log D(x)] + \mathbb{E}_{z\sim p_z(z)}[\log(1-D(G(z)))] .
\end{equation}

理想情况下，生成器能够学习真实数据分布。此时判别器难以区分真实样本和生成样本。GAN 的优点是可以生成较清晰的图像。GAN 的训练也存在问题。常见问题包括模式崩溃、判别器过拟合和样本多样性不足。医学影像数据规模有限，所以普通 GAN 训练可能不稳定。本文采用 StyleGAN2-ADA，是因为该方法在有限数据训练中具有较成熟的工程实现。

\section{StyleGAN2 与 ADA 技术}
StyleGAN 系列方法先把输入潜变量映射到中间潜空间。模型再通过风格调制控制不同层级的图像特征\cite{karras2019stylegan}。这种结构使生成过程更容易控制。StyleGAN2 在此基础上改进了归一化方式。它减少了生成图像中的伪影。它还使用路径长度正则化提升潜空间映射稳定性\cite{karras2020stylegan2}。

StyleGAN2-ADA 的核心是自适应判别器增强。训练数据较少时，判别器可能快速记住训练样本。这样会削弱生成器获得的有效训练信号。ADA 根据训练过程中的统计信号调整增强强度。该机制可以降低判别器过拟合风险\cite{karras2020ada}。本文的 MRI 切片来自有限数量的体数据。虽然切片数量较多，但病例来源仍有限。因此，ADA 有助于提升训练稳定性。

\section{3D Gaussian Splatting 技术}
3D Gaussian Splatting 是一种三维场景表示和渲染方法。它使用许多三维高斯基元表示空间结构。每个高斯包含位置、尺度、方向、不透明度和颜色等参数。与规则体素网格相比，三维高斯可以在连续空间中表达结构。与部分隐式表示相比，3DGS 更便于进行较快的渲染展示\cite{kerbl20233dgs}。

对于第 $i$ 个三维高斯，可将其中心记为 $\mu_i$，协方差矩阵记为 $\Sigma_i$，颜色和不透明度分别记为 $c_i$ 与 $\alpha_i$。其空间影响范围可表示为：

\begin{equation}
G_i(x)=\exp\left(-\frac{1}{2}(x-\mu_i)^{T}\Sigma_i^{-1}(x-\mu_i)\right).
\end{equation}

渲染时，系统把三维高斯投影到目标视角。系统再根据深度、不透明度和颜色合成图像。训练时，模型比较渲染图像和目标图像。模型根据差异更新高斯的位置、尺度、方向、不透明度和颜色。原始 3DGS 多用于多视角照片场景。MRI 切片序列不是普通相机照片。因此，本文需要构建规则相机参数，并把连续切片组织为空间点云和监督图像。

在本项目中，3DGS 主要用于三维表达和平台展示。系统先通过 StyleGAN2 生成或组织连续切片。系统再按照 $z$ 轴顺序堆叠切片。系统把非背景像素转换为带灰度颜色的空间点，并形成 PLY 初始点云。该点云既是 3DGS 训练输入，也是平台进行 2D/3D 对比展示的重要材料。

\section{Gradio 交互平台与工程支撑技术}
Gradio 是一个 Python 工具。它可以把机器学习模型封装为 Web 界面\cite{abid2019gradio}。在本文平台中，Gradio 主要用于组织交互流程。用户可以通过界面填写路径、设置参数、启动任务、查看图像和读取日志。

本文平台以 Python 为主要开发语言。系统使用 PyTorch 支撑 StyleGAN2 训练与推理。系统使用 PIL、NumPy 和 nibabel 处理图像及医学体数据。系统使用 Gradio 构建交互界面。系统还通过子进程调度预处理、训练、生成和点云初始化脚本。该技术组合便于复用开源框架，也便于毕业设计阶段的开发和展示。
